% ============================================================
% SECTION 2: RESEARCH FINDINGS (Slides 5–16)
% ============================================================

\section{Research Findings}

% ----------------------------------------------------------
% Slide 5: The Reality Gap
% ----------------------------------------------------------
\begin{frame}{The Reality Gap}
  \twocol{0.48}{
    \alertmetric{2.5$\times$ Worse}{Real-world vs.\ LRS3 benchmark}

    \vspace{4pt}

    \begin{itemize}
      \item Paper: \wergood{25.4\%} WER (TED talks)
      \item Real-world: \werbad{67.0\%} WER (YouTube)
      \item By WER: only \werbad{11.4\%} ``usable''
    \end{itemize}

    \vspace{6pt}
    {\large\textcolor{ArgosTeal}{\textit{``9 out of 10 segments need verification ---\\or do they?''}}}
  }{0.49}{
    \halfwidthfig{P1_quality_tiers.png}
  }

  \note{
    25.4\% to 67.0\% WER = 2.5x worse. Expected for domain mismatch.
    Only 11.4\% ``usable'' by WER standards.
    But WER dramatically overstates failure --- set up IS metric reveal.
    Hook: ``or do they?'' --- transitions to the IS metric.
  }
\end{frame}

% ----------------------------------------------------------
% Slide 6: WER Is Blind
% ----------------------------------------------------------
\begin{frame}[shrink=5]{WER Is Blind}
  \vspace{2pt}
  {\small Both examples have $\approx$30\% WER, but \alert{opposite meaning preservation}:}

  \vspace{6pt}

  \comparisonbox
    {\faIcon{check-circle}~Intelligible (WER 29\%, IS 4.3)}
    {
      \textbf{Reference:} ``work with the team on this project''\\[2pt]
      \textbf{Hypothesis:} ``work with the team on the project''\\[4pt]
      {\scriptsize Minor word substitution --- meaning fully preserved}
    }
    {\faIcon{times-circle}~Unintelligible (WER 33\%, IS 1.2)}
    {
      \textbf{Reference:} ``admiral mcrae visited the fleet''\\[2pt]
      \textbf{Hypothesis:} ``animal migratory visited the fleet''\\[4pt]
      {\scriptsize Phonetically similar but \alert{completely wrong topic}}
    }

  \vspace{8pt}
  \centering
  {\textcolor{ArgosTeal}{\textbf{Same WER, opposite meaning.}} We built a metric to fix this.}

  \note{
    30 seconds max on this slide.
    Same WER, opposite meaning --- WER cannot distinguish these.
    ``admiral'' and ``animal'' are visually identical on lips.
    This motivates the entire IS framework.
  }
\end{frame}

% ----------------------------------------------------------
% Slide 7: The Intelligibility Score
% ----------------------------------------------------------
\begin{frame}[shrink=20]{The Intelligibility Score (IS)}
  \begin{columns}[T]
    \begin{column}{0.40\textwidth}
      \metricbox{39.9\%}{Properly Captured (IS $\geq$ 3.0)}

      \vspace{3pt}

      \begin{itemize}\small
        \item \keyword{3.5$\times$} more than WER's 11.4\%
        \item 6 deterministic signals:\\
              {\scriptsize Semantic, Phonetic, WER, WWER, NEA, Length}
        \item LLM-distilled design
      \end{itemize}

      \vspace{2pt}

      \begin{block}{\scriptsize Cross-Config (16 configs)}
        {\scriptsize LLM judge: \keyword{r = 0.925}, $\kappa$ = 0.77\\
        Stable: Semantic, Phonetic, NEA\\
        Volatile: WER, Length Ratio}
      \end{block}
    \end{column}
    \begin{column}{0.57\textwidth}
      % IS Tier Distribution Table (Slide 7)
\centering
{\small
\begin{tabular}{@{}lccrr@{}}
  \toprule
  \textbf{Tier} & \textbf{Score} & \textbf{Count} & \textbf{\%} & \\
  \midrule
  \cellcolor{TierExcellentBg}\tierexcellent{Excellent} & 4.0--5.0 & 276 & 18.4\% & \textcolor{TierExcellent}{\rule{1.84cm}{6pt}} \\[1pt]
  \cellcolor{TierGoodBg}\tiergood{Good} & 3.0--3.99 & 321 & 21.4\% & \textcolor{TierGood}{\rule{2.14cm}{6pt}} \\[1pt]
  \cellcolor{TierFairBg}\tierfair{Fair} & 2.0--2.99 & 325 & 21.7\% & \textcolor{TierFair}{\rule{2.17cm}{6pt}} \\[1pt]
  \cellcolor{TierPoorBg}\tierpoor{Poor} & 1.0--1.99 & 336 & 22.4\% & \textcolor{TierPoor}{\rule{2.24cm}{6pt}} \\[1pt]
  \cellcolor{TierFailedBg}\tierfailed{Failed} & 0.0--0.99 & 239 & 16.0\% & \textcolor{TierFailed}{\rule{1.60cm}{6pt}} \\
  \midrule
  \multicolumn{5}{@{}l@{}}{\small\textbf{Properly captured (IS $\geq$ 3.0): \keyword{39.9\%}} --- 3.5$\times$ WER's 11.4\%} \\
  \bottomrule
\end{tabular}}

    \end{column}
  \end{columns}

  \note{
    THIS IS THE MAIN RESEARCH CONTRIBUTION.
    IS reveals 39.9\% properly captured vs 11.4\% by WER --- 3.5x more.
    Claude designed the rubric, selected weights, defined tiers.
    Encoded into deterministic, free, decomposable metrics.
    Cross-config: stable across 16 decode configurations.
    Phonetic preservation is \#1 IS correlate (r=0.943).
  }
\end{frame}

% ----------------------------------------------------------
% Slide 8: Failure Mode Taxonomy
% ----------------------------------------------------------
\begin{frame}[shrink=10]{Failure Mode Taxonomy}
  {\small 10 classified failure modes across \textbf{900 failed segments} (IS $<$ 3.0):}

  \vspace{4pt}

  \begin{center}
  \begin{tikzpicture}
    \pgfplotsset{
      failure bar/.style={
        xbar, bar width=6pt, xmin=0, xmax=18,
        width=0.75\textwidth, height=0.55\textheight,
        xlabel={\small Percentage of failed segments (\%)},
        xlabel style={font=\small},
        ytick=data,
        yticklabel style={font=\scriptsize, anchor=east},
        xticklabel style={font=\scriptsize},
        nodes near coords,
        nodes near coords style={font=\tiny, anchor=west},
        every node near coord/.append style={xshift=1pt},
        enlarge y limits=0.08,
        axis line style={ArgosGray},
        tick style={ArgosGray},
        grid=major, grid style={ArgosLightGray, thin},
      }
    }
    \begin{axis}[failure bar,
      symbolic y coords={
        Over-generation,
        Truncation,
        Empty Output,
        Content Word Errors,
        Entity Destruction,
        High Error Rate,
        Hallucination,
        Accumulated Errors,
        Phonetically Similar,
        Topic Drift
      },
    ]
      \addplot[fill=ArgosRed!70, draw=ArgosRed]
        coordinates {
          (0.1,Over-generation)
          (1.1,Truncation)
          (7.8,Empty Output)
          (10.7,Content Word Errors)
          (12.0,Entity Destruction)
          (12.1,High Error Rate)
          (12.3,Hallucination)
          (12.3,Accumulated Errors)
          (15.7,Phonetically Similar)
          (15.9,Topic Drift)
        };
    \end{axis}
  \end{tikzpicture}
  \end{center}

  \vspace{-2pt}
  \notetext{Each roadmap phase targets specific failure categories. No single fix addresses all modes.}

  \note{
    10 failure modes classified across 900 segments.
    Topic Drift \#1 at 15.9\% --- model drifts to unrelated content.
    Phonetically Similar \#2 at 15.7\% --- inherent visual ambiguity.
    Hallucination at 12.3\% --- fluent but fabricated.
    Key insight: failures are diverse, each phase targets specific categories.
  }
\end{frame}

% ----------------------------------------------------------
% Slide 9: Performance Distribution
% ----------------------------------------------------------
\begin{frame}{Performance Distribution Across 13 Experiments}
  \fullwidthfig{09_boxplot_wwer_all_experiments.png}
                {WWER distribution across 13 decode configurations (107 segments each). Configs D and H are outliers.}

  \note{
    Most segments at 50--80\% WER, long tail of catastrophic failures.
    D (lenpen=-0.5) produces 45\% empties.
    H (lenpen=2.0) produces massive hallucinations (mean WER 540\%).
    The ``sweet spot'' is narrow --- beam=20, lenpen=0--1, optional low-temp sampling.
  }
\end{frame}

% ----------------------------------------------------------
% Slide 10: Why the Gap
% ----------------------------------------------------------
\begin{frame}[shrink=8]{Why the Gap: Root Causes}
  \twocol{0.48}{
    \textbf{Three root causes:}
    \begin{enumerate}
      \item \alert{Domain mismatch}\\
            {\scriptsize TED talks $\neq$ YouTube}
      \item \alert{Hallucination}\\
            {\scriptsize LLM prior overwhelms visual signal}
      \item \alert{Short-segment failure}\\
            {\scriptsize Not enough visual context}
    \end{enumerate}

    \vspace{6pt}

    \begin{exampleblock}{\small Topic Analysis}
      \begin{tabular}{@{}lcc@{}}
        \scriptsize\textbf{Topic} & \scriptsize\textbf{IS} & \scriptsize\textbf{Captured} \\
        \scriptsize Business/Finance & \scriptsize\tiergood{3.08} & \scriptsize 57\% \\
        \scriptsize Education & \scriptsize\tierfair{2.67} & \scriptsize 43\% \\
        \scriptsize DIY/Home & \scriptsize\tierpoor{2.13} & \scriptsize 30\% \\
      \end{tabular}
    \end{exampleblock}

    \vspace{2pt}
    \notetext{Formal vocabulary $\rightarrow$ better predictions}
  }{0.49}{
    \halfwidthfig{01_wer_vs_duration.png}

    \vspace{2pt}
    \begin{center}
    {\small\textbf{Length matters:}}\\[2pt]
    \begin{tabular}{@{}lcc@{}}
      \scriptsize\textbf{Words} & \scriptsize\textbf{WER} & \scriptsize\textbf{Captured} \\
      \scriptsize 5--10 & \scriptsize\werbad{74\%} & \scriptsize 32\% \\
      \scriptsize 20+ & \scriptsize\werneutral{55\%} & \scriptsize 49\% \\
    \end{tabular}
    \end{center}
  }

  \note{
    Short segments catastrophic --- 5--10 words only 32\% captured vs 49\% for 20+.
    Business/Finance best at 57\% (formal vocabulary helps LLM).
    DIY/Home worst at 30\% (casual speech, technical jargon).
    Length is the single strongest predictor of success.
  }
\end{frame}

% ----------------------------------------------------------
% Slide 11: Named Entity Accuracy
% ----------------------------------------------------------
\begin{frame}[shrink=5]{Named Entity Accuracy}
  \twocol{0.46}{
    \alertmetric{38.8\%}{NEA F1 --- Misses 61\% of names/numbers}

    \vspace{4pt}

    \begin{itemize}
      \item Names, numbers, places:\\
            \alert{cannot be recovered from context}
      \item Largest success/failure differentiator:
      \begin{itemize}
        \item Success: \wergood{74.0\%} NEA F1
        \item Failure: \werbad{15.7\%} NEA F1
        \item Gap: \textbf{58.3 pp}
      \end{itemize}
      \item 17.3\% of IS variance\\
            \notetext{(disproportionate to 15\% weight)}
    \end{itemize}
  }{0.51}{
    \halfwidthfig{14_nea_vs_wwer_scatter.png}
  }

  \note{
    NEA F1 38.8\% --- misses 61\% of names and numbers.
    These are the words context can't recover --- ``admiral mcrae'' can't be guessed.
    Largest differentiator: 74\% vs 16\%, gap 58.3pp.
    Contributes 17.3\% of IS variance despite only 15\% weight.
  }
\end{frame}

% ----------------------------------------------------------
% Slide 12: 13 Tuning Experiments
% ----------------------------------------------------------
\begin{frame}[shrink=5]{13 Tuning Experiments}
  \begin{columns}[T]
    \begin{column}{0.48\textwidth}
      \includegraphics[width=\linewidth, height=0.42\textheight, keepaspectratio]{P5_tuning_before_after.png}

      \vspace{2pt}
      \begin{block}{\scriptsize Config J (Best --- Full Dataset)}
        {\scriptsize
        \begin{tabular}{@{}lcc@{}}
          & \textbf{Baseline} & \textbf{Config J} \\
          IS & 2.52 & \keyword{2.60} \\
          Captured & 597 & \keyword{622 (+25)} \\
          Empties & 70 & \keyword{0} \\
        \end{tabular}}
      \end{block}
    \end{column}
    \begin{column}{0.48\textwidth}
      \includegraphics[width=\linewidth, height=0.42\textheight, keepaspectratio]{10_empty_and_hallucination_rates.png}

      \vspace{2pt}
      \begin{alertblock}{\scriptsize The Trade-off}
        {\scriptsize
        \begin{tabular}{@{}lcc@{}}
          & \textbf{Baseline} & \textbf{Config J} \\
          Halluc.\ & 307 & \alert{348 (+41)} \\
          Net IS & --- & +0.08 \\
        \end{tabular}}
      \end{alertblock}
    \end{column}
  \end{columns}

  \note{
    13 experiments: beam, lenpen, sampling, temperature.
    Config J (lenpen=1.0, temp=0.5): IS 2.60 vs 2.52, empties eliminated.
    But hallucinations increase by 41 segments --- fundamental tradeoff.
    Net IS gain only +0.08 --- tuning has reached diminishing returns.
  }
\end{frame}

% ----------------------------------------------------------
% Slide 13: Limits of Tuning
% ----------------------------------------------------------
\begin{frame}{Limits of Tuning}
  \twocol{0.48}{
    {\large\keyword{Tuning = mitigation, not cure}}

    \vspace{4pt}

    \begin{itemize}\small
      \item Config J: eliminates empties,\\doubles hallucinations
      \item Rankings \keyword{stable} (r $>$ 0.92)\\across all 16 configs
      \item ``Hard'' and ``easy'' segments\\stay the same
    \end{itemize}

    \vspace{4pt}

    \begin{alertblock}{\scriptsize Bottleneck: Visual Encoder}
      {\small Three levers remain:}
      \begin{enumerate}\scriptsize
        \item Scale data to 20K--50K segments
        \item Swap to stronger LLM
        \item Smart prompts
      \end{enumerate}
    \end{alertblock}
  }{0.49}{
    \halfwidthfig{P4_lenpen_sensitivity.png}

    \vspace{2pt}
    \notetext{Length penalty vs empty/hallucination trade-off}
  }

  \note{
    Tuning cannot fix the fundamental problem.
    Cross-config proof: segment rankings stable (r>0.92).
    The visual encoder determines which segments succeed or fail.
    Three remaining levers: data, LLM, and prompts.
    This sets up the entire Future Directions section.
  }
\end{frame}

% ----------------------------------------------------------
% Slide 14: Curated Examples
% ----------------------------------------------------------
\begin{frame}[shrink=15]{Curated Examples}
  % Curated Examples Table (Slide 14)
\centering
{\scriptsize
\renewcommand{\arraystretch}{1.3}
\begin{tabular}{@{}p{1.6cm}p{3.8cm}p{4.2cm}rr@{}}
  \toprule
  \textbf{Category} & \textbf{Reference} & \textbf{Hypothesis} & \textbf{WER} & \textbf{IS} \\
  \midrule
  \cellcolor{TierExcellentBg}
  {\tierexcellent{Perfect}} &
  ``\ldots keeping up with what's going on in school so here's the latest taylor and rob\ldots'' &
  33 words about insurance policies, \textbf{every word correct} &
  \wergood{0\%} & \tierexcellent{5.0} \\

  \cellcolor{TierGoodBg}
  {\tiergood{Near-Miss}} &
  ``admiral mcrae visited the fleet'' &
  ``\textbf{animal migratory} visited the fleet'' &
  \werneutral{33\%} & \tiergood{3.0} \\

  \cellcolor{TierFailedBg}
  {\tierfailed{Hallucination}} &
  ``carry strap'' &
  ``I'm so happy to be here with you all today thank you so much for inviting me\ldots'' &
  \werbad{100\%} & \tierfailed{0.7} \\

  \cellcolor{TierFairBg}
  {\tierfair{Tuning Fix}} &
  ``hey it's me just doing the eush\ldots'' &
  \textbf{Baseline}: \textit{(empty)} $\rightarrow$ \textbf{Config J}: ``hi it's me jessica dude who's keeping up\ldots'' &
  \werneutral{73\%} & \tierfair{1.8} \\

  \bottomrule
\end{tabular}}


  \note{
    Walk through each row:
    1. Perfect: 33 words about insurance, WER 0\% --- it works.
    2. Near-miss: ``admiral'' vs ``animal'' --- lip shapes identical.
    3. Hallucination: model fabricates entire narrative. Reads fluently.
    4. Tuning tradeoff: Config J recovers empty output but at cost.
  }
\end{frame}

% ----------------------------------------------------------
% Slide 15: Live Demo
% ----------------------------------------------------------
\begin{frame}{Live Demo}
  \vspace{4pt}
  {\small Sequence: \textcolor{ArgosGreen}{\textbf{Good}} $\rightarrow$ \textcolor{ArgosOrange}{\textbf{Funny}} $\rightarrow$ \textcolor{ArgosRed}{\textbf{Failure}}}

  \vspace{8pt}

  \begin{columns}[c]
    \begin{column}{0.31\textwidth}
      \videocaptioned{\linewidth}%
        {d8BR6hsvzoY_31__2e9546df_thumb.png}%
        {videos/source/d8BR6hsvzoY_31__2e9546df_with_hyp.mp4}%
        {\textcolor{ArgosGreen}{\textbf{WER 0\%}} --- ``buy one get one free''}
    \end{column}
    \begin{column}{0.31\textwidth}
      \videocaptioned{\linewidth}%
        {-POZpyVCN8k_9__c7b26ea8_thumb.png}%
        {videos/source/-POZpyVCN8k_9__c7b26ea8_with_hyp.mp4}%
        {\textcolor{ArgosOrange}{\textbf{WER 33\%}} --- ``admiral'' $\rightarrow$ ``animal''}
    \end{column}
    \begin{column}{0.31\textwidth}
      \videocaptioned{\linewidth}%
        {00MUdHQ7GGY_8__b1480c7a_thumb.png}%
        {videos/source/00MUdHQ7GGY_8__b1480c7a_with_hyp.mp4}%
        {\textcolor{ArgosRed}{\textbf{WER 100\%}} --- hallucinated narrative}
    \end{column}
  \end{columns}

  \vspace{10pt}
  \centering
  \textit{``This is why confidence scoring is the immediate next step.''}

  \note{
    Play videos in order: good, funny, failure.
    1. ``buy one get one free'' --- short, punchy, perfect.
    2. ``admiral mcrae'' to ``animal migratory'' --- audience will laugh.
    3. ``David Irving'' fabrication --- model generates a completely fictional narrative.
    Transition: ``This is why confidence scoring is the immediate next step.''
  }
\end{frame}

% ----------------------------------------------------------
% Slide 16: Research Breadth
% ----------------------------------------------------------
\begin{frame}{Research Breadth: 8 Reports}
  \vspace{2pt}
  \begin{columns}[T]
    \begin{column}{0.48\textwidth}
      \begin{enumerate}
        \item \textbf{Executive Assessment}\\
              {\scriptsize Baseline evaluation, quality tiers, reality gap}
        \item \textbf{Hyperparameter Tuning}\\
              {\scriptsize 13 experiments, parameter sweep}
        \item \textbf{Prompt Engineering}\\
              {\scriptsize Context injection, 7 prompt strategies}
        \item \textbf{Confidence Scoring}\\
              {\scriptsize Beam scores, quality filtering}
      \end{enumerate}
    \end{column}
    \begin{column}{0.48\textwidth}
      \begin{enumerate}
        \setcounter{enumi}{4}
        \item \textbf{N-Best Aggregation}\\
              {\scriptsize ROVER, MBR, hypothesis combination}
        \item \textbf{Fine-Tuning Analysis}\\
              {\scriptsize LoRA r=16/r=64, overfitting analysis}
        \item \textbf{Intelligibility Methodology}\\
              {\scriptsize IS design, 6 signals, tier validation}
        \item \textbf{LLM Upgrade Analysis}\\
              {\scriptsize Model alternatives, data scaling, GER}
      \end{enumerate}
    \end{column}
  \end{columns}

  \vspace{8pt}
  \centering
  \notetext{All reports available as DOCX/PDF with supporting data (CSV, JSON, interactive HTML reports)}

  \note{
    Quick overview --- one sentence per report.
    Each report maps to a backlog mission.
    All available as formatted documents with full data.
    This demonstrates systematic, thorough exploration of the problem space.
  }
\end{frame}
